\clearpage\label{answers}
{\sans\bfseries\LARGE\color{ink} Answers}\par
\footnotesize\linespread{0.94}\selectfont
Try every puzzle first. Where an answer shows a calculation, check the steps, not just the result.

\subhead{Chapter 1}
\textbf{Multix on the 12 units.} No. They went somewhere we did not measure, most likely warmth and sound. ``Destroyed'' claims more than we know. \textbf{Keep box.} Kept: the drop height and the bounce height. Set aside: sound and warmth, not measured.

\textbf{1.1} $50-35=15$.\\
\textbf{1.2} $150\div5=30$~cm. Lost $40-30=10$.\\
\textbf{1.3} (a) S (b) M (c) W (d) S (e) M.\\
\textbf{1.4} Different balls and different drop heights, so the ball could not be the only reason.\\
\textbf{1.5} $28-20=8$. $40-14=26$.\\
\textbf{1.6} $80-56=24$.\\
\textbf{1.7} (b).\\
\textbf{1.8} The empty circle is 12.\\
\textbf{1.9} ``About 12 of the ball's 40 units went into sound and warmth.''\\
\textbf{1.10} Answers vary. Keep the drop height and the floor the same.

\subhead{Chapter 2}
\textbf{Multix on the zero.} No. The motion became warmth in the floor and ball, too small to feel. The zero shows the motion is gone, not the energy. \textbf{Keep box.} Kept: the roll distance for a ramp height. Set aside: the rubbing at each point.

\textbf{2.1} $180\div3=60$~cm.\\
\textbf{2.2} Towel. $150-25=125$~cm.\\
\textbf{2.3} The ramp height and the surface both differed, so the floor and the ramp could not be separated.\\
\textbf{2.4} $40\times7.5=300$~cm.\\
\textbf{2.5} Motion 0, warmth gained 20 (about, in the ledger's units).\\
\textbf{2.6} (a) S (b) M (c) W.\\
\textbf{2.7} A rug rubs more, so the motion becomes warmth sooner.\\
\textbf{2.8} Circles: 20, then 20 as motion, then 0 motion and 20 warmth.\\
\textbf{2.9} ``On the rug, the ball's motion turned into warmth sooner than on the floor.''\\
\textbf{2.10} Answers vary. A higher ramp or a smoother surface roll farther. Keep the ball the same.

\subhead{Chapter 3}
\textbf{Multix on free heat.} No. The warmth came from the motion of the hands, and that came from your body's own energy. Nothing was made from nothing. \textbf{Keep box.} Kept: start and end readings. Set aside: where in the room the warmth went.

\textbf{3.1} $50-37=13$.\\
\textbf{3.2} Open lost 14, lidded lost 8. $14-8=6$ more.\\
\textbf{3.3} (a) motion to warmth (b) stored to motion.\\
\textbf{3.4} The starts differed by 10, so a lid was not the only difference.\\
\textbf{3.5} The air and the table.\\
\textbf{3.6} $14-8=6$.\\
\textbf{3.7} Usually the small cup, since there is less water to hold the warmth. Test both with the same start and lid.\\
\textbf{3.8} Cup circle 15, room circle gains the same 15.\\
\textbf{3.9} ``Warmth left the cup and went into the surroundings.''\\
\textbf{3.10} Answers vary. A lid and a wrap, with readings every 10 minutes.

\subhead{Chapter 4}
\textbf{Multix on saving energy.} No. A lever needs less push but the push is over a longer distance, so the product stays the same or grows. \textbf{Keep box.} Kept: effort times distance. Set aside: friction inside.

\textbf{4.1} $12\div4=3$~cm.\\
\textbf{4.2} $8\times3=24$. $24\div2=12$~cm.\\
\textbf{4.3} Half the pull is 2.5. $50\div2.5=20$~cm.\\
\textbf{4.4} (a) P (b) R (c) L (d) W.\\
\textbf{4.5} $8\times5=40$. $40\div20=2$ coins. The helper pushes 4 times farther.\\
\textbf{4.6} Half the pull needs twice the string, 20~cm.\\
\textbf{4.7} Gap 4. Friction in the wheel.\\
\textbf{4.8} The empty circle is 4.\\
\textbf{4.9} ``A lever trades a smaller push for a longer push.''\\
\textbf{4.10} Answers vary. Measure the pull and the distance, and compare the products.

\subhead{Chapter 5}
\textbf{Multix on the running total.} No. The total lost is $12+8+6=26$. Earlier losses count. \textbf{Keep box.} Kept: the total lost so far. Set aside: what each loss became.

\textbf{5.1} $10$, $8$, $6$. Total 24. $16+24=40$.\\
\textbf{5.2} 14, 24, 32.\\
\textbf{5.3} Bounce 1. The ball has the most energy there, so there is more to lose.\\
\textbf{5.4} $12+8+6+4=30$.\\
\textbf{5.5} 26.\\
\textbf{5.6} $60-26=34$.\\
\textbf{5.7} Warmth and sound.\\
\textbf{5.8} Circles: 40; 28 with 12; then 20 with 8.\\
\textbf{5.9} ``After 4 bounces the ball had 10 units left and 30 had gone into warmth and sound.''\\
\textbf{5.10} Answers vary. You add each deposit and check the jar's total.

\subhead{Chapter 6}
\textbf{Multix on ``destroyed''.} No. The ledger placed the gap at bounce 3 and a chair leg matched it. The energy went into the chair and into sound. \textbf{Keep box.} Kept: where the loss happened. Set aside: the cause of the loss.

\textbf{Spot the cheat.} B. 6 units appeared from nowhere. C has a named gap, friction (warmth).\\
\textbf{6.1} Step 2. $29-24=5$.\\
\textbf{6.2} Yes. Energy appeared from nowhere. The ball may have been pushed down, or the heights misread.\\
\textbf{6.3} No. $3\times10=30$ and $2\times10=20$.\\
\textbf{6.4} No. A 1~cm gap is within reading error.\\
\textbf{6.5} Checking each bounce shows the step where a gap appeared.\\
\textbf{6.6} $0+2+1=3$.\\
\textbf{6.7} (a) M (b) R (c) M.\\
\textbf{6.8} No. The ledger shows where, not why. You would need to watch or repeat the drop.\\
\textbf{6.9} The floors differ, so losses are not comparable.\\
\textbf{6.10} $20-10=10$. $12+8+10=30$.\\
\textbf{6.11} ``4 extra units were lost at bounce 3 and went into the chair and sound.''\\
\textbf{6.12} Check after every step, use the same ball and floor, or name each gap.\\
\textbf{6.13} Calling the whole loss ``warmth'' hides a bounce that lost extra to something else.\\
\textbf{6.14 to 6.16} Answers vary. Check that every circle's number is accounted for.\\
\textbf{Slop check.} (a) You cannot tell: no title, no labels. (b) Many colours, stars, joke text. (c) A title, the heights and the units. (d) The stars and jokes.

\clearpage\label{notes}
{\sans\bfseries\LARGE\color{ink} Words}\par
\small
\begin{multicols}{2}
\word{energy}{what lets something move, warm up or lift; it changes form but the total stays.}
\word{friction}{rubbing between two things that turns motion into warmth.}
\word{ledger}{a list of numbers that adds up to a total.}
\word{lever}{a stiff bar that turns on a point called the pivot.}
\word{model}{a simple stand-in that keeps some things and drops others.}
\word{pulley}{a wheel with a rope that changes the direction or size of a pull.}
\word{ramp}{a slope that trades distance for effort.}
\word{stored energy}{energy a thing has because of where it is, like a lifted ball.}
\end{multicols}

\vspace{0.4em}
{\sans\bfseries\large\color{sky!70!black} Notes for grown-ups}\par
This book teaches that an energy account is trustworthy when what came out adds up to what went in and any gap is named. Each chapter has one activity, one worked example with made-up practice numbers, one trust box, one ``what did we keep?'' prompt and a page of puzzles.

\textbf{Materials and safety.} A bouncy ball, a board and books, a ruler and tape measure, two cups and a lid, a thermometer, coins, a rubber band, string and spools. Use warm tap water only. Keep ramps low and away from stairs. Keep small balls and coins away from young children and pets.

\textbf{Energy units.} Energy is hard to measure at home. The book uses the lift height of the same ball as a stand-in, with 1 unit for each centimetre. The stand-in works only for one ball and one test. Gaps are named, not measured.

\textbf{The practice numbers} were made up so the arithmetic is easy to check. Real results will vary with the ball, the floor and the equipment.

\subhead{Sources}
\begin{itemize}[leftmargin=1.3em,itemsep=0.2em]
\item R. P. Feynman, R. B. Leighton and M. Sands. \emph{The Feynman Lectures on Physics}, Volume I, Chapter~4, ``Conservation of Energy''. For the idea of an energy account.
\item P. G. Hewitt. \emph{Conceptual Physics}. Pearson. For simple machines and for energy changing form.
\item J. R. Taylor. \emph{An Introduction to Error Analysis}, 2nd edition. University Science Books, 1997. For measurement error and repeats.
\end{itemize}

\clearpage
\thispagestyle{blank}
\begin{tikzpicture}[remember picture,overlay]
  \fill[sun!35] ($(current page.south west)+(0.4in,0.4in)$) rectangle ($(current page.north east)-(0.4in,0.4in)$);
  \node[anchor=north,align=center,font=\sffamily\bfseries,text=ink,text width=5.4in] at ($(current page.north)+(0,-1.4in)$) {\Large Next in Science\\[6pt]\fontsize{28}{30}\selectfont Small Models, Big Things};
  \node[align=center,font=\sffamily,text width=4.6in,text=ink] at (current page.center) {A map is smaller than a street. A model is smaller than a room. Book 9 builds one, lists what it keeps and what it drops, and repairs it when it gives a wrong answer.};
  \node at ($(current page.south)+(0,2.6in)$) {\scalebox{3}{\multixbot[ok]}};
  \node[font=\sffamily\bfseries,text=ink] at ($(current page.south)+(0,0.9in)$) {Ages 9 to 11 \quad$\star$\quad Science Book 8};
\end{tikzpicture}
